\section{Evidence behind the recommendations}
\label{sec:discussion-detail}

This appendix gives the evidence and the caveats for the recommendations of Section~\ref{sec:discussion}. Appendix~\ref{sec:deploy} turns them into steps for a verbose generator, and Appendix~\ref{sec:choosing} sets out the three evaluators' measured properties side by side.

\paragraph{The evaluation population.} Where a framework already scores refusals separately, the refusal-pooling result of Table~\ref{tab:4} is a caution about aggregate reporting rather than a defect in the framework. It reached our own pipeline during development, where two implementations of the shared refusal rule disagreed on 365 rows. That is the more realistic failure: the rule exists, and two parts of the pipeline apply different versions of it. We found it before computing any number reported here, and every reported number comes from the single implementation that replaced both.

\paragraph{Claim-level aggregation.} For comparing systems the relevant property is a different one (Appendix~\ref{sec:choosing}). For checking individual answers, claim-level aggregation is a partial remedy. It lifts detection on the verbose generator's NQ cell almost fivefold, from 5.1\% to 25.1\%, and 25.1\% is still a minority of known fabrications, against 61.0\% for the tersest generator under the same aggregate. The gap between generators narrows and does not close, so claim-min reduces the structural bias without removing it. It also costs something the maximum does not. Its absolute scores collapse at that verbosity, from a mean of 0.693 to 0.234 on Natural Questions and 0.424 to 0.076 on HotpotQA, because a minimum over five assertions sits near zero whatever the answer says. Resisting the structural bias and remaining thresholdable are separate properties, and claim-level aggregation buys the first at the cost of the second: a fixed grounding gate carried across to it unchanged would reject almost every verbose answer. Per-assertion thresholding, or weighting assertions by how much of the answer's content they carry, are the obvious next steps and we did not test them.

\paragraph{The deterministic check.} Requiring every value in an answer to appear in the retrieved context flags every falsification we constructed. That part is circular and we state it as such: our second eligibility rule defines the substituted value as absent from the retrieved context, so a containment test must catch it. The informative comparison is the other direction. The three evaluators of the main text reach 89.4\% at best and 5.1\% at worst on the same cases, and a generative fact-checker 97.5\% and 80.6\% (Appendix~\ref{app:minicheck}). These edits are therefore not hard to catch in general. They are hard for an entailment aggregate to catch, and the check is indifferent to answer length where the aggregates are not. The probe does measure one property of the check, its false-positive rate of 0.8\% to 5.2\% on untouched answers, and that rate is not circular.

For a document question answering system, the practical reading of Table~\ref{tab:2} is that a grounding threshold on a verbose generator catches a minority of falsified values: 5.1\% under the common evaluator, and 25.1\% under claim-min, the best of the three on that cell. That is the case for the deterministic check, and two qualifications on it matter. Numbers that assert nothing about the world, such as chunk indices and ordered-list markers, are not expected in the retrieved text, and counting them raises that rate on Claude's HotpotQA answers to about 39\%. The rest are quantities the answer derives instead of copying, which a literal check also flags, so it complements an entailment score rather than replacing one.

Checking symbolic content against the source is an old idea \citep{goodrich-etal-2019, nan-etal-2021, zhao-etal-2020}, and that numbers are a weak point is documented at the level of the representations too \citep{wallace-etal-2019}. We claim no novelty for it. We do claim the measurement that makes the case for it here, against three current evaluators on identical inputs, and note that a practitioner choosing an instrument needs to know which one works. Entailment is separately a loose proxy for attribution: the AIS framework \citep{rashkin-etal-2023} separates support by an identified source from general plausibility.
